% Independent mathematical derivation. Prefix ETC. Requires amsmath,amssymb,hyperref.
\section{Even theta tensors at a fixed negative heat time}
\label{sec:even-theta-tensor-concentration}

\paragraph{Sources and exact scope.}
The kernel and heat sign are those of Brad Rodgers and Terence Tao,
\href{https://arxiv.org/abs/1801.05914v5}{\emph{The De Bruijn--Newman
constant is non-negative}}, original author file
\nolinkurl{fmp-template.tex}, SHA--256
\begin{center}\small
\nolinkurl{25b015c0fb151f1613247d82cf53a6561a320d4dab2bb4f809482e7c477a5817}.
\end{center}
Lines 1--155, including \texttt{phidef} and \texttt{htdef}, were
read directly. No whole-paper reading or application of a
zero-location theorem is asserted. The programme source
\nolinkurl{HEAT_TRANSPORTED_PRODUCT.tex}, HTP1--24, was read
completely, all 539 lines, at SHA--256
\begin{center}\small
\nolinkurl{897bbe51fd757d8c4926aac5a82d771090e1d1dbaed5a9030fcbdbd6e2b74d14}.
\end{center}
In particular HTP19--20, starting at source lines 357 and 380,
give the positivity and moment bounds re-proved below. The
present calculation retains the original kernel, coordinates,
heat parameter and tensor power. It supplies a quantitative
asymptotic for that entire function and makes no conclusion
about the locations of the zeros of $H_0$ or about RH.

\subsection{The original kernel, moments, and strict real maximum}

For $u\ge0$ and $z\in\mathbb C$ retain
\begin{equation}\tag{ETC1}
\begin{gathered}
 \Phi(u)=\sum_{n\ge1}
 (2\pi^2n^4e^{9u}-3\pi n^2e^{5u})e^{-\pi n^2e^{4u}},\\
 H_t(z)=\int_0^\infty e^{tu^2}\Phi(u)\cos(zu)\,du,
 \qquad h=H_0(0),\\
 m_2=\int_0^\infty u^2\Phi(u)\,du,\qquad
 m_4=\int_0^\infty u^4\Phi(u)\,du.
\end{gathered}
\end{equation}
Every original summand of $\Phi$ is positive: its prefactor
is $\pi n^2e^{5u}(2\pi n^2e^{4u}-3)>0$. The $n=1$ term
and the inequalities $n^2-1\ge3(n-1)$, $e^{4u}\ge1$ give
\begin{equation}\tag{ETC2}
\begin{gathered}
 0<\pi(2\pi-3)e^{9u-\pi e^{4u}}
       \le\Phi(u)\le C_\Phi e^{9u-\pi e^{4u}},\\
 C_\Phi=2\pi^2\frac{1+11c_0+11c_0^2+c_0^3}{(1-c_0)^5},
 \qquad c_0=e^{-3\pi}.
\end{gathered}
\end{equation}
Indeed the lower prefactor satisfies
$2\pi e^{4u}-3\ge(2\pi-3)e^{4u}$. For the upper bound
discard the negative term and factor out
$2\pi^2e^{9u-\pi e^{4u}}$. The remaining series is at most
$\sum_{n\ge1}n^4c_0^{n-1}$; differentiating the geometric
series four times with the operator $c\,d/dc$ gives the
stated rational function. On each bounded $u$ interval the
original series converges uniformly, since its absolute
terms are bounded by a fixed multiple of
$(n^4+n^2)e^{-\pi n^2}$. In particular $\Phi$ is continuous.

For $A,B,p\ge0$ all the following integrals are finite,
with the original constants:
\begin{equation}\tag{ETC3}
\begin{split}
 I_p(A,B)&=\int_0^\infty u^pe^{Au^2+Bu}\Phi(u)\,du\\
 &\le\frac{C_\Phi}{2\pi}
 \exp\left(\frac{3A^2}{32\pi}
       +\frac{(B+9+p)^2}{8\pi}-\frac\pi2\right).
\end{split}
\end{equation}
To prove this, use $u^p\le e^{pu}$ and split
$-\pi e^{4u}$ into two quarters and a half. The estimates
$e^{4u}\ge(32/3)u^4$ and $e^{4u}\ge8u^2$ imply
\[
 Au^2-\frac\pi4e^{4u}\le\frac{3A^2}{32\pi},\qquad
 (B+9+p)u-\frac\pi4e^{4u}
          \le\frac{(B+9+p)^2}{8\pi}.
\]
The remaining integral is at most
$e^{-\pi/2}\int_0^\infty e^{-2\pi u}\,du$
because $e^{4u}\ge1+4u$. This proves ETC3.
It follows that $h,m_2,m_4$ are finite and strictly positive.
The bound $|\cos(zu)|\le e^{|\operatorname{Im}z|u}$,
and the corresponding bounds with extra powers of $u$,
justify differentiation under the integral on every compact
set. Thus $H_0$ is entire and even, is real on the real axis,
and
\begin{equation}\tag{ETC4}
 H_0''(0)=-m_2,\qquad H_0^{(4)}(0)=m_4,\qquad
 |H_0(x+iy)|\le I_0(0,|y|)\quad(x,y\in\mathbb R).
\end{equation}

For every real $v\ne0$ there is an open interval of positive
$u$ on which $|\cos(vu)|<1$. Positivity and continuity of
$\Phi$ therefore give
\begin{equation}\tag{ETC5}
 |H_0(v)|\le\int_0^\infty\Phi(u)|\cos(vu)|\,du<h.
\end{equation}
There is also an elementary proof, on this exact kernel, of
\begin{equation}\tag{ETC6}
                         \lim_{|v|\to\infty}H_0(v)=0.
\end{equation}
Given $\eta>0$, choose $U$ with
$\int_U^\infty\Phi<\eta/3$, using ETC3. Uniform continuity
of $\Phi$ on $[0,U]$ supplies a finite step function $s$
with $\int_0^U|\Phi-s|<\eta/3$. If its successive intervals
are $[b_j,b_{j+1}]$ and its values $s_j$, then, for $v\ne0$,
\[
 \left|\int_0^Us(u)\cos(vu)\,du\right|
 \le\frac{2}{|v|}\sum_j|s_j|.
\]
Choose $|v|$ so large that this is below $\eta/3$.
The two approximation errors are bounded by their
$L^1$ norms because $|\cos(vu)|\le1$. Their sum proves
ETC6 without importing a Fourier decay theorem.

\subsection{Actual heat evolution and the full Gaussian contour shift}

Fix $a>0$ and an integer $k\ge1$, and define precisely
\begin{equation}\tag{ETC7}
 T_tf=\sum_{j\ge0}\frac{(-t)^j}{j!}D^{2j}f,
 \quad D=\frac{d}{dz},\qquad
 P_{2k}(-a,z)=T_{-a}(H_0^{2k})(z).
\end{equation}
The series for this actual function will be shown absolutely
convergent; no formal heat operator is substituted for it.
For compatibility with the programme domain,
$H_0^{2k}\in\mathcal E$, where
$\mathcal E=\{f\text{ entire}:\sup_z|f(z)|e^{-\alpha|z|^2}
<\infty\text{ for all }\alpha>0\}$.
In fact $|z|u\le b|z|^2+u^2/(4b)$ gives
$|H_0(z)|\le e^{b|z|^2}I_0(1/(4b),0)$.
Take $b=\alpha/(2k)$ and raise to the exact power $2k$.
This is the same domain and heat sign as HTP2 and HTP5.

Set $\psi(s)=\Phi(|s|)/2$ for $s\in\mathbb R$.
The evenness of the cosine and absolute integrability give
\begin{equation}\tag{ETC8}
\begin{gathered}
 H_0(z)=\int_{\mathbb R}\psi(s)e^{izs}\,ds,\\
 H_0(z)^{2k}=\int_{\mathbb R^{2k}}
       e^{izS}\prod_{j=1}^{2k}\psi(s_j)\,d\boldsymbol s,
 \qquad S=\sum_{j=1}^{2k}s_j.
\end{gathered}
\end{equation}
These are absolutely convergent integrals, since
$|S|\le\sum_j|s_j|$. Differentiation of any finite order
under them is justified, for example, by absorbing
$|S|^N$ into a constant times $e^{|S|}$ and using ETC3.
For the absolute sum of the entire heat series on $|z|\le R$
one has the integrable majorant
\begin{equation}\tag{ETC9}
\begin{split}
 &\int_{\mathbb R^{2k}}e^{R|S|}
       \sum_{j\ge0}\frac{a^jS^{2j}}{j!}
                         \prod_{\ell=1}^{2k}\psi(s_\ell)
                            \,d\boldsymbol s\\
 &\hspace{10mm}=\int_{\mathbb R^{2k}}e^{aS^2+R|S|}
                         \prod_{\ell=1}^{2k}\psi(s_\ell)
                            \,d\boldsymbol s
       \le I_0(2ka,R)^{2k}<\infty.
\end{split}
\end{equation}
Here $S^2\le2k\sum_\ell s_\ell^2$ and
$|S|\le\sum_\ell|s_\ell|$ give the last inequality.
The derivative $D^{2j}e^{izS}=(-1)^jS^{2j}e^{izS}$
now proves, with every exchange absolutely justified,
\begin{equation}\tag{ETC10}
 P_{2k}(-a,z)=\int_{\mathbb R^{2k}}e^{-aS^2+izS}
                    \prod_{\ell=1}^{2k}\psi(s_\ell)
                         \,d\boldsymbol s.
\end{equation}
This also verifies the original negative-time sign.

The real Gaussian identity, with its exact factor, is
\begin{equation}\tag{ETC11}
 \frac{1}{\sqrt{4\pi a}}\int_{\mathbb R}
       e^{-x^2/(4a)}e^{iSx}\,dx=e^{-aS^2}
                    \qquad(S\in\mathbb R).
\end{equation}
One direct proof is to differentiate the left integral in
$S$. Integration by parts, with the Gaussian boundary term
zero, gives $J'(S)=-2aS J(S)$. At $S=0$ its
value is $1$, by squaring the real Gaussian integral and
using polar coordinates. The differential equation gives
ETC11. All differentiations are dominated by a polynomial
times the same Gaussian.
Substituting ETC8 into the following real Gaussian integral
and using ETC11 proves
\begin{equation}\tag{ETC12}
 P_{2k}(-a,z)=\frac{1}{\sqrt{4\pi a}}\int_{\mathbb R}
             e^{-x^2/(4a)}H_0(z+x)^{2k}\,dx.
\end{equation}
The absolute multiple integral before the substitution is
at most $I_0(0,|\operatorname{Im}z|)^{2k}$, since
$|e^{ixS}|=1$ and the Gaussian measure in ETC11 has mass one.
Thus this Fubini interchange is independent of a formal
complex Gaussian manipulation.

We now prove the required contour shift in full. Write
$z=\xi+i\eta$ and put
$F_z(w)=e^{-(z-w)^2/(4a)}H_0(w)^{2k}$, an entire function
of $w$. Formula ETC12 is its integral along
$\mathbb R+i\eta$, oriented with increasing real part,
divided by $\sqrt{4\pi a}$; the real displacement $\xi$
does not change the complete horizontal line. On the closed
strip between heights $0$ and $\eta$, ETC4 gives
$|H_0(w)|\le I_0(0,|\eta|)$. For $|z|\le R$, $M>R$,
and $w=\pm M+iy$ on either vertical edge of that strip,
\begin{equation}\tag{ETC13}
 |F_z(\pm M+iy)|
 \le I_0(0,R)^{2k}
       \exp\left(\frac{R^2}{4a}-\frac{(M-R)^2}{4a}\right),
 \qquad y\text{ between }0\text{ and }\eta.
\end{equation}
Indeed
$\operatorname{Re}(z-w)^2=(\xi\mp M)^2-(\eta-y)^2$,
$|\xi\mp M|\ge M-R$, and $|\eta-y|\le R$.
Each edge has length at most $R$, so both vertical integrals
tend to zero, uniformly on $|z|\le R$ for each fixed $k$.
Both horizontal integrals converge absolutely: on the real
line $|H_0(v)|\le h$, and on the other line its bound is
$I_0(0,R)$, while the real part of the Gaussian supplies
quadratic decay in the horizontal coordinate. Applying
Cauchy's theorem to this rectangle and letting $M$ tend to
infinity, with the same orientation on the two horizontal
lines, proves the entire-function identity
\begin{equation}\tag{ETC14}
 \boxed{\displaystyle
 P_{2k}(-a,z)=\frac{1}{\sqrt{4\pi a}}\int_{\mathbb R}
              e^{-(z-v)^2/(4a)}H_0(v)^{2k}\,dv.}
\end{equation}
For a zero-height strip the equality is immediate. Negative
$\eta$ gives the same equality by reversing the rectangle's
orientation. Thus no sign or half-plane is omitted.

\subsection{Quantitative local concentration and the outer gap}

Define constants from the original moments by
\begin{equation}\tag{ETC15}
 \beta=\frac{m_2}{h}>0,\qquad c=\frac{m_4}{24h}>0,
 \qquad \delta=\min\left\{\beta^{-1/2},
                       \sqrt{\frac{\beta}{4c}}\right\}>0,
 \qquad D_0=\max\left\{c,\frac{\beta^2}{4}\right\}.
\end{equation}
These are fixed constants of $H_0$; $\delta$ is a proved
choice, not an assumed Taylor neighborhood. The inequalities
$\cos x\ge1-x^2/2$ and
$|\cos x-1+x^2/2|\le x^4/24$ give
\begin{equation}\tag{ETC16}
 0\le \frac{H_0(v)}{h}-1+\frac{\beta v^2}{2}\le cv^4
                       \qquad(v\in\mathbb R).
\end{equation}
The first cosine inequality follows from
$1-\cos x=2\sin^2(x/2)\le x^2/2$; the second is Taylor's
integral remainder with the fourth derivative bounded by
$1$. Integrating against the positive original kernel
proves ETC16, with no unknown remainder constant.
For $|v|\le\delta$ it follows that
\begin{equation}\tag{ETC17}
 \frac12\le1-\frac{\beta v^2}{2}\le \frac{H_0(v)}{h}
        \le1-\frac{\beta v^2}{4}\le1,\qquad
 \log\left(\frac{H_0(v)}{h}\right)\le-\frac{\beta v^2}{4},\qquad
 \left|\log\left(\frac{H_0(v)}{h}\right)+\frac{\beta v^2}{2}\right|\le D_0v^4.
\end{equation}
For the last bound write
$e(v)=\frac{H_0(v)}{h}-1+\beta v^2/2\in[0,cv^4]$ and
$y=1-\frac{H_0(v)}{h}=\beta v^2/2-e(v)\in[0,1/2]$.
The convergent real logarithm series yields
\[
 \log\left(\frac{H_0(v)}{h}\right)+\frac{\beta v^2}{2}
 =e(v)-\sum_{n\ge2}\frac{y^n}{n},\qquad
 0\le\sum_{n\ge2}\frac{y^n}{n}
       \le\frac{y^2}{2(1-y)}\le\frac{\beta^2v^4}{4}.
\]
The difference therefore lies between
$-\beta^2v^4/4$ and $cv^4$, proving the maximum in ETC15.
The bound on $\log(H_0(v)/h)$ follows also from $\log(1-y)\le-y$
and $y\ge\beta v^2/4$.

The exact outer gap is
\begin{equation}\tag{ETC18}
 \rho=\sup_{|v|\ge\delta}|\frac{H_0(v)}{h}|,\qquad
 \frac12\le\rho<1,\qquad \lambda=-2\log\rho>0.
\end{equation}
Indeed $\frac{H_0(\delta)}{h}\ge1/2$ by ETC17. By ETC6 there is $M$
so large that $|\frac{H_0(v)}{h}|<1/4$ for $|v|>M$. Hence the
supremum equals a maximum on the compact set
$\delta\le|v|\le M$. Each point of this set has modulus
strictly less than $1$ by ETC5, so its attained maximum is
less than $1$. This proves every assertion in ETC18 and
defines an actual kernel constant, rather than postulating
a uniform outer gap.

For a fixed complex disc $|z|\le R$, $R\ge0$, put
\begin{equation}\tag{ETC19}
\begin{gathered}
 g_z(v)=e^{-v^2/(4a)}\cosh\left(\frac{zv}{2a}\right),\\
 G_R=\frac1{4a}+\frac{R^2}{8a^2}
                      \cosh\left(\frac{R\delta}{2a}\right),
 \qquad
 L_R=\frac{6\sqrt2D_0}{\beta^2}+\frac{\sqrt2G_R}{\beta}.
\end{gathered}
\end{equation}
For real $v$ with $|v|\le\delta$,
\begin{equation}\tag{ETC20}
 |g_z(v)-1|\le G_Rv^2,\qquad
 |\left(\frac{H_0(v)}{h}\right)^{2k}-e^{-\beta kv^2}|
          \le2kD_0v^4e^{-\beta kv^2/2}.
\end{equation}
For the first inequality use $1-e^{-x}\le x$ for $x\ge0$
and
$|\cosh w-1|\le(|w|^2/2)\cosh|w|$, which follows
term by term from the absolutely convergent series.
For the second, the two real exponents
$2k\log\left(\frac{H_0(v)}{h}\right)$ and $-\beta kv^2$ differ by at most
$2kD_0v^4$ and are both at most $-\beta kv^2/2$.
The integral form of the mean value formula for the
exponential proves the second inequality.

Evenness of $H_0^{2k}$ and absolute Gaussian integrability in
ETC14 permit averaging $v$ and $-v$, so that exactly
\begin{equation}\tag{ETC21}
\begin{gathered}
 P_{2k}(-a,z)=\frac{h^{2k}e^{-z^2/(4a)}}{\sqrt{4\pi a}}
                              J_k(z),\\
 J_k(z)=\int_{\mathbb R}g_z(v)\left(\frac{H_0(v)}{h}\right)^{2k}\,dv,
 \qquad M_k=\int_{\mathbb R}e^{-\beta kv^2}\,dv
                    =\sqrt{\frac\pi{\beta k}}.
\end{gathered}
\end{equation}
In particular the even complex weight has no linear term.
On $|v|\le\delta$ the exact decomposition
\[
 g_z(v)\left(\frac{H_0(v)}{h}\right)^{2k}-e^{-\beta kv^2}
 =\bigl(\left(\frac{H_0(v)}{h}\right)^{2k}-e^{-\beta kv^2}\bigr)
       +(g_z(v)-1)\left(\frac{H_0(v)}{h}\right)^{2k}
\]
and ETC17, ETC20 bound the integral of its absolute value
by the whole-real-line integral of
$(2kD_0v^4+G_Rv^2)e^{-\beta kv^2/2}$.
The Gaussian moments
$\int_{\mathbb R}v^2e^{-bv^2}\,dv=\sqrt\pi/(2b^{3/2})$
and
$\int_{\mathbb R}v^4e^{-bv^2}\,dv=3\sqrt\pi/(4b^{5/2})$
follow by differentiating the Gaussian integral in $b>0$;
the derivatives are absolutely dominated on each compact
positive $b$ interval. Substitution of $b=\beta k/2$
and division by $M_k$ give exactly
\begin{equation}\tag{ETC22}
 \frac1{M_k}\left|\int_{-\delta}^{\delta}
    \bigl(g_z(v)\left(\frac{H_0(v)}{h}\right)^{2k}-e^{-\beta kv^2}\bigr)\,dv\right|
                         \le\frac{L_R}{k}.
\end{equation}

The omitted Gaussian tail satisfies
\begin{equation}\tag{ETC23}
 \frac1{M_k}\int_{|v|>\delta}e^{-\beta kv^2}\,dv
                              \le e^{-\beta\delta^2 k}.
\end{equation}
On each half-line set $|v|=\delta+x$, $x\ge0$, and use
$(\delta+x)^2\ge\delta^2+x^2$. The remaining two
half-Gaussian integrals sum to $M_k$.
For the actual outer integral, the series gives
$|\cosh(zv/(2a))|\le\cosh(Rv/(2a))$ for real $v$,
where the latter is even and positive. ETC18 therefore
gives, by completing the square in each of the two real
exponentials making up the hyperbolic cosine,
\begin{equation}\tag{ETC24}
\begin{split}
 \frac1{M_k}\left|\int_{|v|>\delta}g_z(v)\left(\frac{H_0(v)}{h}\right)^{2k}\,dv\right|
 &\le\frac{\rho^{2k}}{M_k}\int_{\mathbb R}
                 e^{-v^2/(4a)}\cosh\left(\frac{Rv}{2a}\right)\,dv\\
 &=2\sqrt{a\beta}\,e^{R^2/(4a)}\sqrt{k}\,\rho^{2k}.
\end{split}
\end{equation}
No decay rate for $H_0(v)$ beyond ETC6 was used; the
Gaussian supplies the absolute outer integrability.

\subsection{Uniform relative error and the exact nonvanishing consequence}

Combining ETC21--24 proves the following quantitative bound
for every $k\ge1$ and every $|z|\le R$:
\begin{equation}\tag{ETC25}
\boxed{\begin{gathered}
 \left|
 \frac{P_{2k}(-a,z)}{
       h^{2k}e^{-z^2/(4a)}\sqrt{h/(4akm_2)}}-1
 \right|\le E_{k,R},\\
 E_{k,R}=\frac{L_R}{k}+e^{-\beta\delta^2k}
          +2\sqrt{a\beta}\,e^{R^2/(4a)}\sqrt{k}\,\rho^{2k}.
\end{gathered}}
\end{equation}
All denominators here are nonzero: $h,a,m_2,k>0$ and the
complex exponential never vanishes. In particular the
factor in the denominator is exactly the factor obtained
from ETC21; its $\pi$ cancels between the two Gaussian
integrals, with no change of the original $a$ or moments.

An explicit relative $O(1/k)$ constant is
\begin{equation}\tag{ETC26}
 \boxed{\displaystyle
 C_R=L_R+\frac1{\mathrm e\beta\delta^2}
       +2\sqrt{a\beta}\,e^{R^2/(4a)}
                         \left(\frac{3}{2\mathrm e\lambda}\right)^{3/2},
 \qquad E_{k,R}\le\frac{C_R}{k}.}
\end{equation}
Here $\mathrm e$ is the exponential constant. For $x>0$,
the maximum of $xe^{-bx}$ is $1/(\mathrm e b)$ for
$b>0$, by direct differentiation. Similarly the maximum
of $x^{3/2}e^{-\lambda x}$ is
$(3/(2\mathrm e\lambda))^{3/2}$. Apply the first with
$b=\beta\delta^2$ to $k e^{-\beta\delta^2k}$ and the
second to $k^{3/2}\rho^{2k}=k^{3/2}e^{-\lambda k}$.
This proves ETC26 for every positive integer $k$.

Consequently the actual even tensor has the locally uniform
complex asymptotic
\begin{equation}\tag{ETC27}
 P_{2k}(-a,z)=h^{2k}e^{-z^2/(4a)}
             \sqrt{\frac{h}{4akm_2}}\,(1+\varepsilon_{k}(z)),
 \qquad \sup_{|z|\le R}|\varepsilon_k(z)|\le\frac{C_R}{k}.
\end{equation}
The statement is uniform in $z$ on each specified disc,
with $a>0$ fixed; it does not assert uniformity as $a$
tends to zero. It gives the explicit lower bound
\begin{equation}\tag{ETC28}
 \inf_{|z|\le R}|P_{2k}(-a,z)|
 \ge h^{2k}\sqrt{\frac{h}{4akm_2}}
              e^{-R^2/(4a)}\left(1-\frac{C_R}{k}\right)>0
                        \qquad(k>C_R).
\end{equation}
Indeed $|e^{-z^2/(4a)}|=e^{-\operatorname{Re}(z^2)/(4a)}
\ge e^{-R^2/(4a)}$ on this disc, and
$|1+\varepsilon_k(z)|\ge1-C_R/k$. Thus every fixed
disc is free of zeros of $P_{2k}(-a,\cdot)$ for all integers
$k>C_R$. This is an unconditional statement about the
actual sequence of functions in ETC7.

In particular, retaining the actual root set of the original
negative-time function, one obtains
\begin{equation}\tag{ETC29}
 \{z:|z|\le R,\ H_{-a}(z)=0,\ P_{2k}(-a,z)=0\}
                      =\varnothing\qquad(k>C_R).
\end{equation}
This follows directly from ETC28 and neither presupposes
nor produces a nonreal root of $H_{-a}$. Any actual fixed
root of $H_{-a}$, real or nonreal, is covered by taking a
disc containing it. The functions $P_{2k}(-a,\cdot)$ are
the heat transforms of the powers $H_0^{2k}$; no argument
here replaces them by the ordinary powers $H_{-a}^{2k}$.
Their quantitatively proved nonvanishing in ETC28 therefore
does not assert nonvanishing of $H_{-a}$ or of $H_0$.
